\expandafter\newcommand\csname captionFigure1\endcsname{\textbf{BioMoAX overview.} (A) Model architecture and five contrastive objectives for learning MOA and compound-identity representations. Colors in the MOA space indicate mechanism groups; positions are schematic. Compound heads use drug measurements only. (B) Retrieval tasks across modalities, batches, cell lines, and perturbation types. Within-cell-line, cross-batch, and cross-cell-line tasks retrieve drugs sharing an MOA. Microscopy images illustrate query and reference inputs.}

\expandafter\newcommand\csname captionFigure2\endcsname{\textbf{Compound identity and MOA retrieval within cell lines.} (A) Same-compound retrieval with disjoint query (20\%) and reference (80\%) crops. Shading shows between-checkpoint SD for BioMoAX and five-split SD for baselines. (B) MOA retrieval with the query compound excluded from the reference library; cross-modal baselines use ridge regression. Error bars show 95\% query-bootstrap CIs; annotations compare BioMoAX with each baseline using Holm-adjusted two-sided Wilcoxon signed-rank tests ($*p<0.05$, $**p<0.01$, $***p<0.001$; ns, not significant, $p\ge0.05$). (C) Image-crop embeddings, with the same annotated crops highlighted across methods. Color denotes MOA, shape distinguishes drugs within each MOA, and gray marks background crops. Two-dimensional maps are fitted independently.}

\expandafter\newcommand\csname captionFigure3\endcsname{\textbf{Cross-batch MOA retrieval and representation structure.} (A) Images of 115 cpg0030 drugs query the cpg0012 reference library for other drugs sharing an MOA; identical-compound matches are absent. Error bars show checkpoint SD for BioMoAX and 95\% query-bootstrap intervals for baselines. Stars show Holm-adjusted two-sided Wilcoxon comparisons with each baseline; thresholds follow Fig.~2. (B) Batch mixing for 32 compounds shared by the two sources; higher normalized scores indicate greater mixing. (C) Two-dimensional distributions of image-crop representations, colored by batch (top) and MOA (bottom). Each point represents a crop; coordinates are identical between the two rows for each method. Gray indicates other or missing MOA annotations.}

\expandafter\newcommand\csname captionFigure4\endcsname{\textbf{Mechanism-related retrieval across cell lines.} (A) Bidirectional U2OS--A549 retrieval in four modality combinations; arrows indicate query-to-reference direction. (B) HT29 images, excluded from image training, query the pooled U2OS/A549 image (I2I) or L1000 expression (I2E) library. In A and B, hits are reference drugs sharing an MOA with the query. Error bars show 95\% query-bootstrap CIs; significance notation follows Fig.~2. (C) Drug-level image representations colored by cell line (top) and MOA (bottom), with shapes denoting cell line. Each point represents an aggregated drug profile; both rows use identical coordinates within each method. Highlighted MOA groups are selected by proximity in the BioMoAX map; gray denotes other drug profiles. Lines connect highlighted profiles sharing an MOA.}

\expandafter\newcommand\csname captionFigure5\endcsname{\textbf{Drug retrieval from CRISPR perturbations.} (A) Retrieval of 1,038 reference drugs using 29 held-out A549 gene perturbations. Positives comprise target-annotated drugs and one-hop MOA neighbors; shading shows 95\% query-bootstrap CIs. (B,C) Selected co-retrieval examples from individual checkpoints: KDR/DDR2--linifanib (B) and ALDH2/AKR1B1--disulfiram (C). Solid pink edges mark the reference-annotated KDR--linifanib and ALDH2--disulfiram pairs; dashed ochre edges mark additional co-retrieval without direct-target annotation. Boxes are enlarged in the adjacent insets; numbers link nodes to image examples.}
